\subsection{The ordinary two-prime class and the necessary balanced correction}

\begin{statement}{A.168}{Proved}
{A nonzero ordinary ideal class at every higher level}
Retain the fixed tower and the distinguished degree-one primes
\(P_{2,0,n}\) of A.166.  Then
\[
             [P_{2,0,n}]_{(3)}\ne0
                  \quad\hbox{in }\operatorname{Cl}(K_n)\{3\}
                  \qquad(n\geq3).
\tag{A.168.1}
\]
In particular no integral element of \(K_n\), for \(n\geq3\),
has rational norm \(2\) or \(-2\).

Let \(\mathcal F_{3,S}\) be the completed \(S\)-unit norm fibre
\[
 \mathcal F_{3,S}=\{x\in\widehat E_{3,3}:
                                  N_{K_3/\Q}x=[-2]\}.
\]
Its actual valuation map at the two quadratic changed-prime branches
has the following proved image containment:
\[
 \begin{gathered}
 \nu:\mathcal F_{3,S}\longrightarrow\Z_3^2,
       \qquad x\longmapsto(c_{11},c_{59}),\\
 \nu(\mathcal F_{3,S})\subset
       \{(b,c)\in\Z_3^2:c-4b\equiv3\pmod9\}.
 \end{gathered}
\tag{A.168.2}
\]
At either changed prime the full divisor vector is
\((-2c_q,c_q)\), with the actual residue-degree weights \((1,2)\).
The same constraint holds for the stable changed-prime coordinates
of every unrestricted completed coherent norm family over \([-2]\),
through A.166.  In particular at least one of these coordinates
has three-adic valuation at most one.

For the balanced classes of A.155 one has the sharper bound
\[
                         3^{n-2}e_n=0\qquad(n\geq2).
\tag{A.168.3}
\]
The quotient map from the ordinary class group to
\(C_n^{\rm bal}\) takes the nonzero class in A.168.1 to \(e_n\).
Neither A.168.1 nor A.168.3 decides whether that image is zero.
The balanced class \(e_3\), and the coherent norm fibre, remain
unresolved by these assertions.
\end{statement}

\begin{proof}
Suppose the three-primary projection of \([P_{2,0,3}]\) were zero.
Its order \(h\) in the finite ordinary class group would then be
prime to three.  There would be an actual element
\(\gamma\in K_3^\times\) with
\[
             \operatorname{div}(\gamma)=hP_{2,0,3}.
\tag{A.168.4}
\]
Its rational norm is \(\pm2^h\).  The element
\(x=\gamma^{h^{-1}}\) is well defined in the pro-three completion,
using the inverse of \(h\) in \(\Z_3\); its completed norm is
\([-2]\), since the sign has trivial pro-three class.
It belongs to \(\widehat E_{3,3}\) and its divisor is exactly
\(P_{2,0,3}\).

The distinguished local factors all have residue degree one, so
\[
 N_{3,2}P_{2,0,3}=P_{2,0,2},\qquad
 N_{3,1}P_{2,0,3}=P_{2,0,1}.
\tag{A.168.5}
\]
These are ideal norms, with their literal residue-degree factors.
Since \(\operatorname{div}(\beta_2)=P_{2,0,2}\) and
\(\operatorname{div}(\beta_1)=P_{2,0,1}\), the actual elements
\[
   N_{3,2}(\gamma)/\beta_2^h\in\mathcal O_{K_2}^{\times},
   \qquad
   N_{3,1}(\gamma)/\beta_1^h\in\mathcal O_{K_1}^{\times}
\tag{A.168.6}
\]
are units.  Taking their completed \(h\)-th roots and using the
proved pro-three unit basis \(U\) in \(K_1\), the first seed is
therefore exactly \(\beta_1U^z\) for some \(z\in\Z_3\).
The second seed is a completed \(S\)-unit preimage, and it is
itself a completed global norm from \(K_3\), hence a local norm
at every place.

Apply the exact finite-height locus A.165.4 with \(a=b=c=0\).
Its two changed-prime equations require respectively
\[
                        z\equiv6\pmod9,
                  \qquad z\equiv3\pmod9.
\tag{A.168.7}
\]
They are incompatible.  This proves A.168.1 for \(n=3\).
For every higher \(n\), ideal norm sends
\(P_{2,0,n}\) to \(P_{2,0,3}\), again with weight one.
It respects primary projections of the finite class groups,
so it proves the assertion for every \(n\geq3\).
There is only one prime of residue degree one above two:
the unramified polynomial \(X^{3^n}-a\) has the unique simple
root one modulo two.  An integral element of rational norm
\(\pm2\) would therefore generate \(P_{2,0,n}\), contradicting
the proved nonzero primary class.

For A.168.2, take \(x\in\mathcal F_{3,S}\) and set
\(b_i=N_{3,i}x\).  The complete first fibre in A.165.2 writes
\(b_1=\mathfrak b(z,a,b,c)\) uniquely.  The second seed is an
\(S\)-unit preimage passing every local \(K_3/K_2\) norm test.
Thus A.165.4 gives
\(b\equiv z+3\) and \(c\equiv4z+6\pmod9\).
Eliminating \(z\) gives
\(c-4b\equiv3\pmod9\).
At each changed prime all extensions after \(K_1\) are totally
ramified on the two fixed branches.  Their ideal-norm weights are
one, so the first-seed quadratic valuations \(b,c\) are exactly
the coordinates \(c_{11},c_{59}\) of \(x\).  The rational norm
condition gives the accompanying linear coordinates \(-2c_q\).
This proves the stated image containment.  A.166 puts every
unrestricted coherent family in this fixed-support fibre and
identifies its stable coordinates, proving the further assertion.
If both coordinates were divisible by nine, their displayed
combination could not be congruent to three modulo nine.

Finally extend the principal ideal \((\beta_2)=P_{2,0,2}\)
from \(K_2\) to \(K_n\).  The entire two-adic tower is unramified,
so its extension divisor is the sum, with coefficient one, of
the primes above \(P_{2,0,2}\).  In the precise balanced quotient
A.155.2 each such prime has class \(f_P e_n\).  The sum of
these absolute residue degrees is \([K_n:K_2]=3^{n-2}\), because
the old prime has residue degree one and all ramification indices
are one.  The class of the principal extension divisor is zero,
which proves A.168.3.  This bounds the order of the element;
it does not bound the precision needed to detect it in its ambient
finite group, as already explained in A.155.
\end{proof}

The finite congruence elimination, compatible two-prime cycle weights,
and balanced extension-degree sums are checked in
\path{code/two_inert_raw_two_class_certificate.py}, with retained output.
